\subsection{Cambio de anillos}
\label{subsection:0.6.2-fr}

Cuando un grupo aditivo $M$ está provisto de varias estructuras de módulo con
respecto a anillos $A$, $B$, \ldots, en lugar de decir que $M$ es plano
como $A$-módulo, $B$-módulo, \ldots, diremos a veces también que $M$
es $A$-\emph{plano}, $B$-\emph{plano}, \ldots

\begin{env}[6.2.1]
\label{0.6.2.1-fr}
Sean $A$ y $B$ dos anillos, $M$ un $A$-módulo, $N$ un
$(A,B)$-bimódulo. Si $M$ es plano y si $N$ es $B$-plano, entonces
$M\otimes_A N$ es $B$-plano. En particular, si $M$ y $N$ son dos
$A$-módulos planos, $M\otimes_A N$ es un $A$-módulo plano. Si $B$ es una
$A$-álgebra y si $M$ es
\oldpage[0\textsubscript{I}]{56}
un $A$-módulo plano, el $B$-módulo $M_{(B)}=M\otimes_A B$ es plano. Finalmente, si
$B$ es una $A$-álgebra que es plana como $A$-módulo, y si $N$ es
un $B$-módulo plano, entonces $N$ es también $A$-plano.
\end{env}

\begin{env}[6.2.2]
\label{0.6.2.2-fr}
Sean $A$ un anillo, $B$ una $A$-álgebra plana como $A$-módulo.
Sean $M$, $N$ dos $A$-módulos, tales que $M$ admite una presentación
finita; entonces el homomorfismo canónico
\[
  \operatorname{Hom}_A(M,N)\otimes_A B
  \longrightarrow
  \operatorname{Hom}_B(M\otimes_A B,N\otimes_A B)
  \tag{6.2.2.1}
\]
(que transforma $u\otimes b$ en el homomorfismo
$m\otimes b'\to u(m)\otimes b'b$) es un isomorfismo.
\end{env}

\begin{env}[6.2.3]
\label{0.6.2.3-fr}
Sea $(A_\lambda,\varphi_{\mu\lambda})$ un sistema inductivo filtrante
de anillos; sea $A=\varinjlim A_\lambda$. Sea, por otra parte, para cada
$\lambda$, $M_\lambda$ un $A_\lambda$-módulo y para $\lambda\leq\mu$ sea
$\theta_{\mu\lambda}:M_\lambda\to M_\mu$ un
$\varphi_{\mu\lambda}$-homomorfismo, tal que
$(M_\lambda,\theta_{\mu\lambda})$ es un sistema inductivo;
$M=\varinjlim M_\lambda$ es entonces un $A$-módulo. Dicho esto, si para todo
$\lambda$, $M_\lambda$ es un $A_\lambda$-módulo \emph{plano}, entonces $M$ es
un $A$-módulo \emph{plano}. En efecto, sea $\mathfrak{J}$ un ideal \emph{de
tipo finito} de $A$; por definición del límite inductivo, existe un
índice $\lambda$ y un ideal $\mathfrak{J}_\lambda$ de $A_\lambda$ tales que
$\mathfrak{J}=\mathfrak{J}_\lambda A$. Si se pone
$\mathfrak{J}'_\mu=\mathfrak{J}_\lambda A_\mu$ para $\mu\geq\lambda$, se tiene
también $\mathfrak{J}=\varinjlim\mathfrak{J}'_\mu$ (donde $\mu$ recorre los
índices $\geq\lambda$), de donde (siendo el funtor $\varinjlim$ exacto y
conmutando con el producto tensorial)
\[
  M\otimes_A\mathfrak{J}
  =\varinjlim(M_\mu\otimes_{A_\mu}\mathfrak{J}'_\mu)
  =\varinjlim\mathfrak{J}'_\mu M_\mu
  =\mathfrak{J}M.
\]
\end{env}

